\chapter{Szybki przewodnik po Mathematice}

\begin{summarybox}
Ten dodatek jest praktycznym skrótem poleceń Mathematiki potrzebnych w całych
notatkach. Nie zastępuje dokumentacji programu, ale zbiera najczęściej używane
konstrukcje: definicje funkcji, listy, wykresy, upraszczanie, rozwiązywanie
równań, całkowanie, równania różniczkowe, animacje i eksport wyników.
\end{summarybox}

\section{Definiowanie funkcji}

Funkcję definiujemy najczęściej z użyciem wzorca podkreślenia:
\begin{verbatim}
f[x_] := x^2 + 1
f[3]
\end{verbatim}
Znak \texttt{:=} oznacza definicję opóźnioną: prawa strona jest obliczana dopiero wtedy, gdy funkcja zostanie wywołana. Znak \texttt{=} wykonuje obliczenie od razu.

Dla funkcji wielu zmiennych:
\begin{verbatim}
psi[x_, t_] := Exp[-(x - t)^2]
\end{verbatim}
Warto pilnować, aby argumenty funkcji zawsze miały podkreślenia w definicji. Bez nich Mathematica definiuje wartość tylko dla konkretnego symbolu.

\section{Zmienne lokalne i globalne}

Do zmiennych lokalnych używamy \texttt{Module}:
\begin{verbatim}
g[a_] := Module[{x}, x = a + 1; x^2]
\end{verbatim}
To chroni przed przypadkowym użyciem globalnej wartości \texttt{x}. W większych obliczeniach jest to bardzo ważne, bo wiele błędów wynika z pozostawionych wcześniej definicji.

Do czyszczenia symboli używamy:
\begin{verbatim}
Clear[x, y, psi]
ClearAll["Global`*"]
\end{verbatim}
Druga komenda czyści wszystkie symbole w bieżącym kontekście i jest przydatna na początku notebooka.

\section{Listy, tabele i mapowanie}

Lista w Mathematice ma postać:
\begin{verbatim}
{1, 2, 3, 4}
\end{verbatim}
Tabele generujemy przez \texttt{Table}:
\begin{verbatim}
Table[n^2, {n, 1, 10}]
Table[Sin[n x], {n, 1, 5}]
\end{verbatim}
Mapowanie funkcji po liście zapisujemy jako:
\begin{verbatim}
f /@ {1, 2, 3}
\end{verbatim}
Albo pełniej:
\begin{verbatim}
Map[f, {1, 2, 3}]
\end{verbatim}

\section{Funkcje czyste}

Funkcja czysta to szybka funkcja anonimowa. Na przykład:
\begin{verbatim}
(#^2 + 1) & /@ {1, 2, 3}
(#1 + #2) &[3, 4]
\end{verbatim}
Symbol \texttt{\#} oznacza argument, a \texttt{\&} kończy definicję funkcji.

\section{\texttt{Plot}, \texttt{ListPlot}, \texttt{DensityPlot}, \texttt{DensityPlot3D}}

Podstawowy wykres funkcji jednej zmiennej:
\begin{verbatim}
Plot[Sin[x], {x, 0, 2 Pi}]
\end{verbatim}
Kilka funkcji na jednym wykresie:
\begin{verbatim}
Plot[{Sin[x], Cos[x]}, {x, 0, 2 Pi}, PlotLegends -> Automatic]
\end{verbatim}
Wykres punktowy:
\begin{verbatim}
ListPlot[Table[{x, Sin[x]}, {x, 0, 10, 0.2}]]
\end{verbatim}
Mapa gęstości 2D:
\begin{verbatim}
DensityPlot[Sin[x y], {x, -3, 3}, {y, -3, 3}]
\end{verbatim}
Trójwymiarowa mapa gęstości:
\begin{verbatim}
DensityPlot3D[Exp[-x^2 - y^2 - z^2], {x, -2, 2}, {y, -2, 2}, {z, -2, 2}]
\end{verbatim}

\section{\texttt{Evaluate} w wykresach}

Przy wykresach rozwiązań symbolicznych albo numerycznych często warto wymusić wcześniejsze podstawienie wyniku:
\begin{verbatim}
Plot[Evaluate[y[t] /. sol], {t, 0, 5}]
\end{verbatim}
Bez \texttt{Evaluate} Mathematica czasem próbuje analizować wyrażenie zbyt późno, co może spowolnić wykres albo wygenerować ostrzeżenia.

\section{\texttt{Simplify}, \texttt{FullSimplify}, \texttt{Assumptions}}

Upraszczanie bez założeń bywa niewystarczające:
\begin{verbatim}
Simplify[Sqrt[a^2]]
\end{verbatim}
Mathematica nie wie, czy \(a\) jest dodatnie. Dlatego używamy założeń:
\begin{verbatim}
Simplify[Sqrt[a^2], Assumptions -> a > 0]
\end{verbatim}
Dla bardziej agresywnego upraszczania:
\begin{verbatim}
FullSimplify[expr, Assumptions -> {a > 0, n \[Element] Integers}]
\end{verbatim}

Można także użyć konstrukcji \texttt{Assuming}:
\begin{verbatim}
Assuming[a > 0, FullSimplify[Sqrt[a^2]]]
\end{verbatim}

\section{\texttt{Solve}, \texttt{NSolve}, \texttt{FindRoot}}

Równania symboliczne rozwiązujemy przez:
\begin{verbatim}
Solve[x^2 - 1 == 0, x]
\end{verbatim}
Numeryczne rozwiązywanie wielomianów:
\begin{verbatim}
NSolve[x^5 - x + 1 == 0, x]
\end{verbatim}
Równania transcendentalne często wymagają \texttt{FindRoot}:
\begin{verbatim}
FindRoot[Cos[x] == x, {x, 1}]
\end{verbatim}
W \texttt{FindRoot} punkt startowy ma znaczenie. Różne punkty startowe mogą prowadzić do różnych pierwiastków.

\section{\texttt{NIntegrate}}

Całki numeryczne liczymy przez:
\begin{verbatim}
NIntegrate[Abs[psi[x]]^2, {x, xmin, xmax}]
\end{verbatim}
To jest szczególnie przydatne przy sprawdzaniu normalizacji. Jeśli \texttt{NIntegrate} zgłasza problemy, trzeba sprawdzić osobliwości funkcji, zakres całkowania i założenia numeryczne.

\section{\texttt{NDSolve} i \texttt{InterpolatingFunction}}

Równania różniczkowe numerycznie rozwiązujemy przez \texttt{NDSolve}. Przykład równania zwyczajnego:
\begin{verbatim}
sol = NDSolve[{y'[t] == -y[t], y[0] == 1}, y, {t, 0, 5}]
Plot[Evaluate[y[t] /. sol], {t, 0, 5}]
\end{verbatim}
Wynik \texttt{NDSolve} często zawiera obiekt \texttt{InterpolatingFunction}. Oznacza to, że rozwiązanie jest znane numerycznie jako interpolacja na określonym zakresie. Nie należy używać go poza zakresem, dla którego zostało policzone.

\section{\texttt{Chop}, \texttt{WorkingPrecision} i drobne błędy numeryczne}

W obliczeniach numerycznych mogą pojawiać się bardzo małe liczby, które są tylko szumem numerycznym:
\begin{verbatim}
Chop[{1, 10^-16, -10^-14}]
\end{verbatim}
Jeżeli problem jest czuły numerycznie, można zwiększyć precyzję:
\begin{verbatim}
N[expr, 30]
FindRoot[eqn, {x, x0}, WorkingPrecision -> 40]
\end{verbatim}
Większa precyzja nie naprawia złego modelu, ale pomaga przy trudnych równaniach transcendentalnych.

\section{\texttt{Manipulate} i \texttt{Animate}}

Interaktywny suwak:
\begin{verbatim}
Manipulate[Plot[Sin[a x], {x, 0, 2 Pi}], {a, 1, 10}]
\end{verbatim}
Animacja po czasie:
\begin{verbatim}
Animate[Plot[Sin[x - t], {x, 0, 2 Pi}], {t, 0, 2 Pi}]
\end{verbatim}
W raportach końcowych animacja powinna być uzupełniona statycznymi wykresami i komentarzem fizycznym.

\section{Eksport wykresów i animacji}

Wykres można zapisać do pliku:
\begin{verbatim}
plt = Plot[Sin[x], {x, 0, 2 Pi}];
Export["wykres.pdf", plt]
Export["wykres.png", plt, ImageResolution -> 300]
\end{verbatim}
Do dokumentów najlepiej używać PDF dla wykresów wektorowych i PNG dla map gęstości albo obrazów rastrowych.

Animację można wyeksportować na przykład jako GIF:
\begin{verbatim}
Export["animacja.gif", frames]
\end{verbatim}
W raporcie często wystarczy kilka wybranych klatek animacji i opis, co zmienia się w czasie.

\section{Typowe komunikaty błędów}

Nie należy ignorować czerwonych komunikatów. Najczęstsze problemy to:
\begin{itemize}
    \item brak wartości numerycznej funkcji w \texttt{Plot} albo \texttt{NIntegrate},
    \item zły zakres zmiennej,
    \item osobliwość w całce,
    \item brak warunku początkowego albo brzegowego w \texttt{NDSolve},
    \item próba użycia rozwiązania numerycznego poza zakresem interpolacji.
\end{itemize}

\begin{summarybox}
Najważniejsza zasada pracy w Mathematice: nie ufać samemu obrazkowi. Trzeba kontrolować założenia, zakresy, jednostki, normalizację, komunikaty błędów i sens fizyczny wyniku.
\end{summarybox}
